\section{Introduction}

Sound design is a central discipline in music production. Whether crafting the lead timbre of a synthesizer patch, shaping the tone of a bass line, or sculpting textural pads, producers routinely spend hours adjusting oscillator waveforms, filter cutoffs, envelope shapes, and effect parameters to achieve a desired sonic character. When a musician hears an inspiring sound in an existing recording and wishes to recreate it, the standard workflow is laborious: listen carefully, form a mental model of the signal's spectral and temporal properties, and then manually navigate the high-dimensional parameter space of a synthesizer until the output sounds acceptably close. This trial-and-error process---sometimes called \emph{sound matching} or \emph{preset reverse-engineering}---demands deep expertise and considerable patience, placing it out of reach for many aspiring producers.

Contemporary music technology has made significant progress in extracting \emph{symbolic} information from audio. Source separation systems such as Spleeter \citep{hennequin2020spleeter} and its successors decompose a polyphonic mix into isolated stems (vocals, drums, bass, other), while multi-instrument transcription models like MT3 \citep{gardner2022mt3} can identify which notes are played and by which general instrument class. These tools answer the questions \emph{what notes are being played?} and \emph{what broad category of instrument is this?}---but they stop short of answering the more granular question: \emph{what synthesis parameters would reproduce this particular timbre?} A transcription system may recognise that a passage contains a ``synthesizer lead,'' yet it cannot tell the user which oscillator mix, filter resonance, or envelope decay would recreate that specific sound.

This gap between instrument \emph{classification} and instrument \emph{recreation} has attracted attention in the audio machine-learning community. Engel et al.\ introduced Differentiable Digital Signal Processing (DDSP), demonstrating that embedding traditional signal-processing modules---oscillator banks, filters, reverbs---inside a differentiable computation graph enables gradient-based optimisation of synthesis parameters from audio \citep{engel2020ddsp}. Barkan et al.\ tackled the inverse problem directly with InverSynth, training deep networks to predict synthesizer configurations from spectrograms \citep{barkan2019inversynth}. Despite these advances, practical systems that accept an arbitrary audio recording and return a playable synthesizer patch remain elusive: differentiable synthesizers constrain the timbral palette to sounds they can represent, while learned models require large paired datasets and generalise poorly to unseen synth architectures.

In this paper we present \textsc{Instrumental}, a system that addresses the inverse synthesis problem through a hybrid optimisation pipeline. Given a target audio signal, \textsc{Instrumental} automatically recovers the parameters of a differentiable subtractive synthesizer that, when rendered at the detected pitch, produces audio perceptually matching the input. The pipeline proceeds in three stages: (1)~spectral initialisation, which analyses the target spectrum to produce an informed starting point in parameter space; (2)~global search via Covariance Matrix Adaptation Evolution Strategy (CMA-ES), a derivative-free optimiser that explores the multimodal loss landscape without requiring gradient information; and (3)~local refinement via Adam gradient descent through the fully differentiable synthesizer, exploiting the smooth local geometry to converge on a precise solution. A composite perceptual loss---combining mel-scaled multi-resolution STFT, spectral centroid distance, and MFCC divergence---guides all three stages.

Our contributions are as follows:
\begin{enumerate}
    \item An end-to-end pipeline from audio to playable synthesizer patch, combining spectral analysis, evolutionary search, and gradient-based refinement.
    \item A systematic comparison of optimisation strategies---including CMA-ES, gradient descent, SPSA, multi-start search, and hybrid approaches---for navigating the non-convex parameter landscape of a 28-parameter differentiable synthesizer.
    \item An empirical study of eight hypotheses for improving convergence and timbral fidelity, spanning architectural extensions (FM synthesis, Chebyshev waveshaping, parametric EQ), alternative loss functions (CLAP embeddings, CDPAM perceptual distance, temporal fine-structure loss), and optimiser variants (SPSA fine-tuning, multi-start CMA-ES).
\end{enumerate}
